\documentclass[11pt,a4paper]{article}
\usepackage[margin=25mm]{geometry}
\usepackage{amsmath,amssymb,amsthm,booktabs,hyperref}
\setlength{\emergencystretch}{2em}
\newtheorem{lemma}{Lemma}
\newtheorem{theorem}[lemma]{Theorem}
\newcommand{\gen}[1]{\langle#1\rangle}
\title{Small primitive attachments through a decorated tail}
\author{Human Compression project: bounded L044 proof component}
\date{}
\begin{document}\maketitle
This note proves the attachment exclusion in Proposition~5.11 of
Cameron--Fernandes--Leemans--Mixer, \emph{Highest rank of a polytope
for $A_n$}, arXiv:1605.09173v1, physical page~26. Its sixteen small
markings are supplied by the separately accepted exact Table~2
classification. The source's alternating-attachment idea and its
attribution to Fernandes--Leemans--Mixer are retained. The kernel-orbit
argument below treats the extra central decorations explicitly; it
does not assert IP for a projected tail. No whole-original or
new-question claim is made.

\begin{lemma}[Decorated alternating attachment]\label{lem:decorated}
Let $W,Y,R$ be disjoint finite sets, $Y=\{u,v,w,z\}$, and $c\in W$.
Let $T$ preserve all three sets and be transitive on $W$.
Suppose $K=\ker(T\to S_Y)$ has a $W$-orbit $O$ containing $c$
with $|O|\geq2$. Let $a,b$ be involutions preserving $W\cup Y$ and $R$, with
\[
 a|_{W\cup Y}=(c\ u)(v\ w),\qquad
 b|_{W\cup Y}=(u\ v)(w\ z).
\]
If the action of $K$ on $R$ centralizes $a|_R$, then
$\gen{T,a,b}$ contains the full alternating group on $W\cup Y$,
acting trivially on $R$.
\end{lemma}
\begin{proof}
For every $j\in O$, conjugation by an element of $K$ supplies
\[
                    a_j=(j\ u)(v\ w)\,a|_R.
\]
Products $a_j a_l$ for distinct $j,l$ cancel the decoration and
give pure three-cycles through $u$. They generate
$A_{O\cup\{u\}}$, acting trivially elsewhere. Its conjugate by
$b$ is $A_{O\cup\{v\}}$. Alternating support gluing on their
intersection $O$, of size at least two, gives
$A_{O\cup\{u,v\}}$. Conjugating this by $a$ and gluing gives
$A_{O\cup\{u,v,w\}}$; conjugating that by $b$ and gluing gives
$A_{O\cup Y}$. All these are pure supported alternating groups:
conjugating a permutation supported in $W\cup Y$ cannot introduce
support in the invariant set $R$.

The normal subgroup $K$ has its $W$-orbits permuted transitively
by $T$. Conjugates of $A_{O\cup Y}$ therefore give
$A_{O'\cup Y}$ for every such orbit $O'$. Their common set $Y$
has four points, so repeated alternating support gluing gives
$A_{W\cup Y}$ with trivial action on $R$. The gluing identity
is $(p\ q\ x)^{-1}(p\ q\ y)=(p\ x\ y)$ in right-action
notation, for distinct $p,q$ in the common support. Thus no
classification of $T$ or IP assumption on it is needed.
\end{proof}

\begin{lemma}[The twelve literal attachment templates]\label{lem:templates}
Take a marking $H=A_X$ from rows~4--16 of source Table~2, in
either its printed order or its full reverse. Write its rank
as $d$, $a=g_{d-1}$, $b=g_{d-2}$ and
$H_-=\gen{g_0,\ldots,g_{d-2}}$. Either $H_-$ fixes no point
of $X$, or it fixes exactly one point $c$ and the following
properties hold. There are four distinct points $u,v,w,z$,
all different from $c$, such that on $\{c,u,v,w,z\}$
\[
                   a=(c\ u)(v\ w),\qquad b=(u\ v)(w\ z).
\]
Put $Y=\{u,v,w,z\}$, $R=X\setminus(\{c\}\cup Y)$, and
$C=C_{S_{X\setminus\{c\}}}(H_-)$. Both $a,b$ preserve $R$;
$C$ preserves $Y,R$ and has order at most two. The group induced
on $R$ by $a,b,C$ has order at most two and is even. The action
of $C$ on $Y$ is trivial when $|X|=11$. Finally
$|\gen{a,b}|$ is ten or twenty, hence below sixty.
\end{lemma}
\begin{proof}
This is a finite check of the literal drawings, independent of
the marking census. The mandatory \path{source_table2.json}
is the exact previously matched edge transcription.
\texttt{templates.g} checks each of the 26 orientations. It
computes the entire $H_-$ fixed set and the entire centralizer
$C$, sets
\[
              u=c^a,\quad v=u^b,\quad w=v^a,\quad z=w^b,
\]
and verifies the displayed restrictions, all invariant sets,
the auxiliary action, its parity and the last-pair order.
There are fourteen orientations without an endpoint and twelve
with the stated template. Their exact vertex lists and complete
centralizer elements are in the fresh receipt.

The positive orientations are rows~4 and its reverse; rows~5,
6,7,10; row~11 and its reverse; row~12; row~14; and row~16 and
its reverse. Only rows~5,6,7 have a possibly nontrivial auxiliary
group on $R$. In rows~6 and~7 its generator is the even
involution $(4\ 8)(5\ 9)$ in the literal upper-then-lower vertex
numbering. This decoration is retained, not discarded by an
unproved projection argument. The whole template program and
data are required finite premises; the source is not used as
an assertion that these tests pass.
\end{proof}

\begin{theorem}[Primitive prefix exclusion]\label{thm:prefix}
Let $G=\gen{\rho_0,\ldots,\rho_{r-1}}\leq A_n$ be transitive,
with a connected string C-marking and all maximal omissions
intransitive, where $n\geq12$. Suppose
$H=\gen{\rho_0,\ldots,\rho_{d-1}}$ has nonempty support $X$ and acts
primitively and transitively there, with $m=|X|<12$ and
$d\geq(m-1)/2$. Then its coloured graph is row~1,2 or~3 of
source Table~2, or the appropriate full reverse. The dual
suffix assertion follows by reversing the whole marking.
\end{theorem}
\begin{proof}
The prefix inherits IP, evenness and connectedness. An omission
of one of its generators cannot remain transitive on $X$:
otherwise the corresponding global omission connects every
pair in $X$, so adding the omitted permutation, supported on
$X$, cannot change its orbits. This contradicts global
transitivity. The accepted Table~2 theorem therefore applies.

Suppose its row is among 4--16, so $H=A_X$ and $m=9,10$ or~11.
Since $X$ is proper and $G$ is transitive, $d<r$.
Any permutation commuting with $H$ fixes $X$ pointwise. Indeed
it preserves the nontrivial $H$-orbit $X$, and the centralizer
of the natural alternating action is trivial: the point
stabilizer fixes only that point. Thus every $\rho_j$ with
$j\geq d+1$ fixes $X$ pointwise. The only possible exit from
$X$ is $x=\rho_d$, which centralizes $H_-$. If $t\in X$ is
sent outside, then for every $h\in H_-$,
$(t^h)^x=(t^x)^h=t^x$, so $t^h=t$. The no-endpoint templates
are impossible. In every remaining case the only exit point
is the unique $c$, and $x$ interchanges it with an outside point.

Let $Y,R$ be its template and put
$W=(\Omega\setminus X)\cup\{c\}$,
$T=\gen{\rho_d,\ldots,\rho_{r-1}}$.
The nonfixed set of $H_-$ is $X\setminus\{c\}$, so $x$
preserves it and restricts there to an element of $C$.
Consequently $T$ preserves $W,Y,R$, has image of order at most
two on $Y$, and its action on $R$ belongs to the same even
group of order at most two as $a|_R,b|_R$. In particular it
centralizes $a|_R$.

The tail $T$ is transitive on $W$. To see this, use connectivity
of the full coloured permutation graph. All old prefix edges
are inside $X$, and a tail edge can leave $X$ only at $c$.
Every excursion through $X$ of a walk between points of $W$
can therefore be deleted at $c$; the remaining walk uses only
tail colours in $W$.

Let $K=\ker(T\to S_Y)$. Its normal orbits on $W$ have equal
size and there are at most two. For $m=9$, $|W|\geq4$; for
$m=10$, $|W|\geq3$, and a partition into two equal orbits
cannot have singleton parts. For $m=11$, the template action
on $Y$ is trivial, so $K=T$ and $|W|\geq2$. Thus in every
case a $K$-orbit containing $c$ has at least two points.
Lemma~\ref{lem:decorated} yields
\[
 A_{W\cup Y}\ \leq\ J:=\gen{\rho_{d-2},\ldots,\rho_{r-1}},
\]
with trivial action on $R$. Intersecting with $H=A_X$ gives
$A_{\{c\}\cup Y}\cong A_5\leq H\cap J$. But full marked IP
requires $H\cap J=\gen{a,b}$, of order at most twenty. This
contradiction excludes every row~4--16.
\end{proof}

\paragraph{The source's intermediate groups.}
The argument also justifies the group shapes used in its
case summary. All actions on $R$ in the templates are even,
so, since $G\leq A_n$, every element of $J$ induces an even
permutation on $W\cup Y$. We proved that the pure alternating
factor is contained in $J$. Hence $J=A_{W\cup Y}$ if the
auxiliary action is trivial, and otherwise
$J=A_{W\cup Y}\times C_2$. This is an actual direct product
on disjoint supports. The latter possibility occurs only for
rows~5,6,7. The support size is $n-m+5$, giving the source's
$A_{n-4},A_{n-5},A_{n-6}$ factors respectively. This argument
retains every decoration and does not infer a marked isomorphism
from an abstract group order.

\paragraph{Scope and dependencies.}
The complete Table~2 packet is a paid premise. The whole new
finite dependency is \texttt{templates.g}, \texttt{check.py}
and the exact source graph data; GAP permutation arithmetic
is retained. No subgroup census, source extraction or Lean
rebuild is repeated. The source's subsequent supported-interval
classification and whole L044 remain separate obligations.
\end{document}
